\documentclass{article}

% Idioma e codificação para português brasileiro
\usepackage[brazil]{babel}
\usepackage[utf8]{inputenc}
\usepackage[T1]{fontenc}

% Margens e tamanho do papel
\usepackage[a4paper,top=2cm,bottom=2cm,left=3cm,right=3cm,marginparwidth=1.75cm]{geometry}

% Pacotes úteis
\usepackage{amsmath}
\usepackage{graphicx}
\usepackage{booktabs}
\usepackage[colorlinks=true, allcolors=blue]{hyperref}

\newcommand{\p}{\rule[-0.8mm]{20mm}{3.5mm}}

\title{Caracterização de Padrões de Depreciação Automotiva sob a Transição Energética: Uma Abordagem Baseada em Representação Espacial de Séries Temporais}
\author{Gabriel M. Brum}
\date{2026}

\begin{document}
\maketitle

\begin{abstract}
Este trabalho investiga as dinâmicas temporais de depreciação de veículos automotores a partir do conjunto de dados fipeX , com foco no impacto da inserção de veículos elétricos (EVs) no mercado brasileiro frente aos modelos convencionais a combustão interna . Metodologicamente, o estudo aborda três questões centrais: (i) a verificação de assimetrias estatísticas nas taxas de depreciação entre elétricos e a combustão de mesma faixa de valor; (ii) a análise de co-movimento e efeitos de transbordamento decorrentes de guerras de preços e marcas entrantes ; e (iii) a avaliação comparativa entre a clusterização em séries temporais 1D tradicionais e representações bidimensionais geradas via \textit{Gramian Angular Field} (GAF) e \textit{Markov Transition Field} (MTF). A proposta visa unir a caracterização econômica de novos ativos à exploração de representações espaciais como estratégia de transformação de atributos em tarefas de mineração de dados não supervisionadas.
\end{abstract}

\section{Introdução}

O mercado automotivo global e brasileiro atravessa uma transição tecnológica impulsionada pela entrada acelerada de Veículos Elétricos (EVs) \cite{alvarez2024modelagem}. Diferente dos veículos a combustão interna (ICE), cujo comportamento de depreciação histórica é consolidado e linearizado pela idade do bem, os modelos eletrificados introduzem variáveis de incerteza operacional: degradação acelerada de baterias, rápida obsolescência tecnológica, políticas tarifárias voláteis e guerras de preços \cite{gautam2024depreciation, li2026global}.

Compreender e segmentar essas trajetórias temporais não é apenas uma demanda mercadológica de precificação e análise de risco para seguradoras e instituições financeiras; constitui um problema clássico de \textit{Data Mining} voltado à mineração de séries temporais (\textit{time-series mining}) e descoberta de padrões em dados ordenados. Métodos convencionais frequentemente achatam as curvas históricas ou dependem de distâncias euclidianas 1D que falham em capturar dinâmicas não lineares, flutuações sazonais e correlações defasadas entre segmentos concorrentes \cite{schloter2022empirical}.

A representação convencional de séries temporais como vetores unidimensionais impõe limitações analíticas severas ao processo de mineração de dados, restringindo os algoritmos de agrupamento a métricas de distância ponto a ponto altamente suscetíveis a ruídos, variações de escala e desalinhamentos temporais. Em dados ordenados complexos, como trajetórias de precificação de veículos, o tratamento vetorial tradicional desconsidera interações de vizinhança e estruturas não lineares inerentes às curvas de depreciação. Essa restrição metodológica fundamenta a necessidade de investigar transformações de atributos mais robustas, avaliando se representações bidimensionais geradas por \textit{Gramian Angular Fields} (GAF) ou \textit{Markov Transition Fields} (MTF) são capazes de codificar correlações cruzadas e dependências dinâmicas que escapam às distâncias euclidianas clássicas \cite{wu2023imaging, ulyanin2019feature}.

Sob o ponto de vista da modelagem estatística e do aprendizado de máquina, o estudo se justifica ao contrapor a eficácia de espaços latentes extraídos dessas imagens frente a métodos tradicionais de redução de dimensionalidade e seleção de atributos. A aplicação de matrizes GAF/MTF permite transpor o problema de dados tabulares temporais para o domínio de operadores de extração de padrões espaciais, abrindo caminho para testar se técnicas de aprendizado não supervisionado se beneficiam da preservação topológica dessas transformações \cite{jueco2024performance}. Com isso, o trabalho estabelece um teste empírico controlado para mensurar se a complexidade adicional dessa projeção espacial se traduz, de fato, em \textit{clusters} mais coesos e discriminantes do que os obtidos via abordagens vetoriais convencionais.

Por fim, o trabalho preenche uma lacuna relevante na análise de fenômenos de mercado ao estruturar a descoberta de padrões de desvalorização e possíveis efeitos de transbordamento decorrentes da chegada dos veículos elétricos e das montadoras chinesas \cite{li2025rise, abrardi2026market}. Ao analisar o co-movimento de preços entre matrizes energéticas concorrentes, a pesquisa vai além do agrupamento meramente exploratório e fornece um protocolo experimental rigoroso baseado em técnicas consolidadas de pré-processamento, mitigação do impacto de \textit{outliers} e validação de modelos. Dessa forma, a investigação entrega simultaneamente uma contribuição teórica sobre representação de dados ordenados na mineração de dados e uma ferramenta descritiva aplicada para precificação e gestão de risco sobre ativos automotivos.

\section{Objetivos}

\begin{itemize}
      \item \textbf{Mapear e quantificar as taxas de depreciação dos veículos elétricos:} Mensurar a velocidade e amplitude da perda de valor dos modelos eletrificados em comparação direta a veículos a combustão pareados por porte, categoria e faixa de preço de lançamento \cite{schloter2022empirical}.
      
      \item \textbf{Avaliar o efeito de transbordamento (\textit{spillover}) de marcas entrantes:} Analisar correlações temporais e possíveis quebras estruturais de preços em segmentos tradicionais após a introdução e guerra tarifária promovida por montadoras chinesas no mercado nacional \cite{li2026global, li2025rise}.

      \item \textbf{Implementar e validar o \textit{reshaping} de séries temporais em imagens 2D:} Transformar as trajetórias unidimensionais de preços de cada veículo em representações matriciais utilizando algoritmos como \textit{Gramian Angular Field} (GAF) ou \textit{Markov Transition Field} (MTF) \cite{wu2023imaging, jueco2024performance}.

      \item \textbf{Identificar e caracterizar arquétipos de desvalorização veicular:} Aplicar algoritmos de agrupamento (\textit{clustering}) sobre representações vetoriais clássicas e sobre os espaços latentes gerados a partir das imagens, identificando perfis comportamentais típicos (alta retenção, desvalorização acelerada, valorização tardia).

      \item \textbf{Comparar empiricamente a eficácia das representações:} Avaliar quantitativamente, por meio de métricas de qualidade de agrupamento (como coeficiente de Silhueta e índice Davies-Bouldin), se a transformação das séries para imagem supera as abordagens vetoriais unidimensionais tradicionais.
\end{itemize}

\section{Atividades}

\begin{itemize}
      \item \textbf{Atividade 1:} Coleta, limpeza e pré-processamento dos dados, tratamento de inconsistências/valores ausentes e estruturação temporal dos registros da base fipeX por veículo.
      
      \item \textbf{Atividade 2:} Análise exploratória e pareamento comparativo visualização das trajetórias de depreciação, cálculo de estatísticas descritivas e pareamento dos modelos elétricos com equivalentes a combustão para identificação de padrões preliminares.
      
      \item \textbf{Atividade 3:} Transformação e representação espacial (\textit{reshaping}), realizando a codificação das séries temporais 1D em representações matriciais 2D utilizando as técnicas de GAF e MTF.
      
      \item \textbf{Atividade 4:} Modelagem de agrupamento e extração de padrões aplicação e ajuste dos algoritmos de \textit{clustering} (K-Means e abordagens baseadas em densidade) tanto sobre os dados vetoriais clássicos quanto sobre as representações geradas por imagem.

      \item \textbf{Atividade 5:} Avaliação experimental, interpretação e escrita final com o cálculo das métricas de validação de agrupamento, análise dos arquétipos comportamentais descobertos e consolidação dos resultados no artigo final.
\end{itemize}
  
\section{Cronograma}

\begin{table}[h]
\centering
\scalebox{0.82}{
\begin{tabular}{||c||c|c|c|c|c|c||}
\hline\hline
\multicolumn{1}{||c||}{Atividades} & Semana 1 & Semana 2 & Semana 3 & Semana 4 & Semana 5 & Semana 6 \\ \hline\hline
 1 & \p &    &    &    &    &    \\ \hline
 2 &    & \p &    &    &    &    \\ \hline
 3 &    &    & \p &    &    &    \\ \hline
 4 &    &    &    & \p & \p &    \\ \hline 
 5 &    &    &    &    &    & \p \\ \hline \hline
\end{tabular}}
\caption{Cronograma de execução das atividades em semanas.}
\label{tab:cronograma}
\end{table}

\bibliographystyle{plain}
\bibliography{references}

\end{document}